\scanpage{9}
\textcircled{2.4} \unsure{Une} donnée de descente sur $(\unsure{\mathscr X|\mathcal X|\mathfrak X}'\mathbin{-}\mathscr X)/\bar k$ de $\bar k$ vers $k$, revient à la donnée d’un scindage
\diagram{(20)\quad\begin{tikzcd}[ampersand replacement=\&]
\unsure{\mathcal G|\mathscr G}'\arrow[r,dash]\&\Gamma\arrow[l,bend left=38]
\end{tikzcd}}
permettant d’écrire
\[
\mathcal G'\simeq\Gamma\mathbin{\underset{1/2}{\cdot}}\deleted{\unsure{\mathcal E}}\unsure{\mathfrak G|G|\mathscr G}'.\tag*{(21)}
\]
Plus généralement, si $k\subset k_1\subset\bar k$ correspond à $\Gamma_1\subset\Gamma$ ss-gpe fermé, une \inserted{donnée de} descente de $\bar k$ à $k_1$ revient à un scindage partiel
\diagram{(22)\quad\begin{tikzcd}[ampersand replacement=\&,row sep=small]
\mathcal G'\arrow[r,dash]\&\Gamma\\
\&\Gamma_1\arrow[ul]\arrow[u,phantom,"\cup" description]
\end{tikzcd}}
et quitte à introduire les variantes
% slanted
\margin{oubli à \unsure{réparer|vérifier} ds rédaction en forme!}
\[
\left|\begin{aligned}
\unsure{\mathcal E|\mathscr E}_1&=\mathrm{Aut}_{k_1}(\widetilde{\mathscr X}\mathbin{-}\mathscr X)\unsure{,}=\text{l’image inverse de }\Gamma_1\subset\Gamma\\
\mathcal E'_1&=\ldotp\ldotp\ldotp\ldotp=\mathcal E'\cap\mathcal E_1\qquad\text{par }\begin{tikzcd}[ampersand replacement=\&,column sep=small]
\mathcal E\arrow[r,dash]\arrow[rr,bend right=23]\&\mathcal G\arrow[r]\&\Gamma
\end{tikzcd}\\
\mathcal G'_1&=\mathrm{Aut}_{k_1}(\mathscr X'\mathbin{-}\mathscr X)=\text{l’image inverse de }\Gamma_1\subset\Gamma\\
&\qquad\text{par }\mathcal G'\mathbin{-}\mathcal G\mathbin{-}\Gamma
\end{aligned}\right.\tag*{(23)}
\]
il revient au $=$ de se donner un iso. compatible avec la structure d’ext
\[
\mathcal G'_1\simeq\Gamma_1\mathbin{\underset{1/2}{\cdot}}\mathfrak G'
\]
% archive sheet number: 8
